\section{Problem Setup and Backbone Bank}
\label{sec:setup}

The audit evaluates eight frozen STGNN checkpoints under three
deployment-realistic shifts. Table~\ref{tab:audit_design} summarizes each
shift axis, its metric, the protocol control, the main validity threat, and
the appendix that addresses it.

\begin{table}[H]
\centering
\caption{\textbf{Audit design at a glance.} Three deployment-realistic shifts evaluated
under a unified frozen-checkpoint protocol.}
\label{tab:audit_design}
\vskip 0.05in
\small
\begin{tabular}{p{2.3cm}p{2.6cm}p{2.6cm}p{2.8cm}p{1.5cm}}
\toprule
\textbf{Shift axis} & \textbf{Metric} & \textbf{Control} & \textbf{Validity threat} & \textbf{App.\ check} \\
\midrule
Sensor dropout &
  $\DeltaSensor$: MAE delta (norm.) &
  MAE on kept sensors only; primary mask &
  Random vs.\ structured dropout pattern &
  \ref{sec:appendix_structured_dropout}, \ref{sec:appendix_multimask} \\
Cross-city transfer &
  X-city norm.\ MAE &
  Matched structural access (source adj) &
  Subset sampling; padding inflation &
  \ref{sec:appendix_subset_sensitivity}, \ref{sec:appendix_no_pad_sensitivity} \\
Temporal subgroup &
  $\DeltaTemp$: weekend$-$weekday (norm.; mph comparison) &
  Within-split stratification; no new data &
  Metric-dependent sign reversal &
  \ref{sec:appendix_temporal_absolute}, \ref{sec:appendix_temporal_cov} \\
\bottomrule
\end{tabular}
\end{table}

\paragraph{How to read the audit.}
Each column in the main tables is an estimand, not a single leaderboard score.
Clean MAE identifies the checkpoint being audited. $\DeltaSensor$ measures
incremental damage from silent zero-masking on the source city. Cross-city MAE
measures zero-shot transfer under matched structural access, and $\DeltaTemp$
measures a weekday/weekend subgroup gap inside the held-out source-city test
split. We therefore compare models within a shift axis and treat changes in
rank across axes as the object of study. This convention also avoids turning
the masked-MAE convention used by newer adaptive-relation models into a false
clean-performance comparison with the canonical bank.

\subsection{Backbone Architectures}
\label{subsec:backbones}

The canonical bank covers five architectures from graph-free to fixed-graph to
adaptive-graph.

\textbf{DCRNN}~\citep{li2018dcrnn} uses a bidirectional diffusion GRU on a
fixed road-distance adjacency $\Aadj$ (truncated at $K$ hops in both
directions), fixed at training and inference.
\textbf{Graph WaveNet (GWNet)}~\citep{wu2019gwnet} adds a learned adaptive
adjacency
\begin{equation}
  \Aadapt = \text{softmax}(\mathbf{E}_1 \mathbf{E}_2^\top),
  \label{eq:adaptive_adj}
\end{equation}
with $\mathbf{E}_1, \mathbf{E}_2 \in \mathbb{R}^{N \times d_e}$ trained end-to-end and tied to the $N$ training nodes.
\textbf{STGCN}~\citep{yu2018stgcn} alternates gated temporal convolutions with
truncated-Chebyshev spectral graph convolutions over the same fixed road
adjacency.
\textbf{GRU} is a graph-free recurrent baseline that treats the $N$ sensors as
a single feature vector and uses no adjacency.
\textbf{MTGNNLight} is a lightweight MTGNN variant~\citep{wu2020mtgnn} that
learns a sparse adaptive graph via pairwise cosine similarity on a single
embedding $\mathbf{E}\in\mathbb{R}^{N\times d}$ and keeps only top-$k$
neighbors per node. This directly contrasts with GWNet's dense all-pairs
softmax and isolates whether dropout fragility depends on adjacency density or
on topology learning per se (P14).

These five share a unified config (hidden dim 64, seq/pred 12, Adam $1{\times}10^{-3}$,
weight decay $10^{-4}$, cosine schedule, 100 epochs, patience 15). To rule out
training-artifact explanations of GWNet's fragility, we additionally train
\textbf{GWNet[CM]} (patience 30) and \textbf{GWNet[Strong]} (P17; hidden 128,
adaptive embed 20, patience 30); both are frozen at shift-eval time.

\subsection{Backbone Bank Protocol}
\label{subsec:bank}

We train each backbone once on METR-LA~\citep{li2018dcrnn} using the
chronological 70/10/20 split and three fixed seeds. The resulting checkpoints
are frozen before any shift evaluation: no model is re-trained, fine-tuned, or
selected separately for sensor dropout, cross-city transfer, or temporal
subgroup analysis.
An additional four-backbone few-shot probe relaxes this constraint to test
limited target-data adaptation (Appendix~\ref{sec:appendix_fewshot}).

Table~\ref{tab:main} reports the clean METR-LA MAE together with the three
shift metrics. The clean column is included only to anchor the starting point
of each checkpoint. Our claims are not based on clean ranking: a backbone can
be useful for diagnosing a failure mechanism even when it is not the strongest
source-city forecaster. For the same reason, GWNet[CM] and GWNet[Strong] are
treated as control checkpoints for the GWNet finding, not as additional entries
in the main bank.

\subsection{Extended adaptive-relation models}
\label{subsec:extended_models}

To probe generality, we add three architectures using the same METR-LA split
and frozen-checkpoint rule:
\textbf{AGCRN}~\citep{bai2020agcrn} (node-adaptive GRU, no temporal
covariates); \textbf{MegaCRN}~\citep{jiang2023megacrn} (meta-graph memory
with TOD/DOW inputs); \textbf{STAEformer}~\citep{liu2023staeformer}
(Transformer with spatial-temporal adaptive embeddings encoding temporal
covariates). These models use their standard masked-MAE convention, while the
canonical bank uses plain MAE; Table~\ref{tab:extended} reports clean values
only to identify checkpoints, not to rank clean performance across families.

\subsection{Distribution shifts}
\label{subsec:shifts}

\paragraph{Sensor dropout.}
At inference, a random 30\% of sensor channels are zero-masked. We do not
re-instantiate the 207-node model; it sees a 207-dimensional input with zeroed
channels, and MAE is computed on the kept sensors only so clean and masked
evaluations are directly comparable. Tables~\ref{tab:main}--\ref{tab:extended}
report the pre-registered primary-mask aggregate; Appendix~\ref{sec:appendix_multimask}
reports 5-mask aggregates to separate mask-sampling from model-weight variance.
GRU is the main mask-sampling caveat. Structured dropout patterns
(hub and contiguous) are in Appendix~\ref{sec:appendix_structured_dropout}.
This protocol isolates silent failure; the LOCF comparison in
Appendix~\ref{sec:appendix_imputation} shows that effective imputation can
largely remove the zero-mask gap.
Supplementary probes add stale hold, noise, bias, and intermittent missingness
(Appendix~\ref{sec:appendix_corruption_suite}).

\paragraph{Cross-city transfer.}
The frozen METR-LA model is evaluated zero-shot on
PEMS-BAY~\citep{li2018dcrnn} (Bay Area, 325 sensors, 52,116 timesteps).
A second pair (PEMS04 source, 307 sensors; zero-shot on METR-LA and PEMS-BAY)
is in Appendix~\ref{sec:appendix_pems04} for cross-city generality (P18).
``Zero-shot'' means no re-training or fine-tuning. Because PEMS-BAY has 325
sensors and the METR-LA checkpoints have 207 input channels, we subsample 207
PEMS-BAY sensors (5 subsets); rankings remain stable at 150-node coverage
(Appendix~\ref{sec:appendix_subset_sensitivity}).

\textbf{Matched structural access (canonical comparison).}
The main cross-city comparison uses source-only adjacency: every backbone
uses the METR-LA adjacency at inference, removing the asymmetry that
adaptive-graph models cannot access target topology by construction.
Under this constraint, any observed advantage is architectural, not an
artifact of privileged topology access
(target-adj results: Appendix~\ref{sec:appendix_source_adj}).
MAE is normalized with source-city statistics for cross-architecture
comparability.

\paragraph{Temporal subgroup shift.}
METR-LA spans weekdays and weekends with different traffic patterns; the
model is trained on both. We stratify the test set by day type---weekday
(5,104 windows) and weekend (1,728 windows)---and report
$\DeltaTemp = \mathrm{MAE}_{\text{weekend}} - \mathrm{MAE}_{\text{weekday}}$.
This is a subgroup performance gap within the existing test split, not
deployment to a new time window. In addition to the normalized gaps, we report
the extended models' masked-mph gaps on the same checkpoints and windows
(Appendix~\ref{sec:appendix_temporal_absolute}). The latter exclude zero-valued
targets, whereas normalized MAE includes all targets; the comparison changes
both scaling and the target mask.

\paragraph{Protocol summary.}
{\small Train on METR-LA (207 sensors, 5-min, 70/10/20 split); 3 seeds,
100 epochs, patience 15; frozen at evaluation. \emph{Sensor:} 30\% silent
primary zero-mask for Tables~\ref{tab:main}--\ref{tab:extended}, MAE on kept sensors.
\emph{City:} zero-shot to
PEMS-BAY (325$\to$207 subsampled, 5 subsets), source-only adjacency
(matched-structural-access; padding sensitivity in
Appendix~\ref{sec:appendix_no_pad_sensitivity}). \emph{Temporal:} weekday
(5{,}104) vs.\ weekend (1{,}728) test windows. \emph{Aggregation:}
mean$\,\pm\,$population std over 3 seeds unless a table explicitly reports
multi-mask summaries;
METR-LA train-statistics normalization for the canonical bank,
model-specific scaler for extended models
(\S\ref{subsec:extended_models}).}
